\documentclass[10pt,a4paper]{amsart}
\usepackage[margin=19mm]{geometry}
\usepackage{amsmath,amssymb,mathtools}
\usepackage[scr=rsfs,cal=euler]{mathalfa}
\usepackage{xfrac}
\usepackage[unicode,hidelinks,bookmarks=false]{hyperref}
\newcommand{\ie}{i.e.\ }
\newcommand{\cf}{cf.\ }
\newcommand{\resp}{respectively\ }
\newcommand{\Subsection}{\subsection}
\providecommand{\nobreakdash}{\nobreak\hbox{-}}
\setlength{\emergencystretch}{2em}
\multlinegap0pt
\usepackage{bm}
\usepackage{bbm}
\usepackage{dsfont}
\usepackage{xspace}
\usepackage{cancel}
\usepackage{mathtools}
\usepackage{amsthm}
\makeatletter
\@ifundefined{theorem}{\newtheorem{theorem}{Theorem}}{}
\@ifundefined{lemma}{\newtheorem{lemma}[theorem]{Lemma}}{}
\makeatother
\makeatletter
\expandafter\ifx\csname innerproof\endcsname\relax
\newenvironment{innerproof}[1]{
\begin{proof}[Proof of #1]\renewcommand{\qedsymbol}{\ensuremath{\textit{(End of proof of #1)\quad $\blacksquare$}}}}{\end{proof}}
\fi
\renewcommand{\leq}{\leqslant}
\makeatother
\title[From the Finite to the Infinite: Sharper Asymptotic Bounds o]{From the Finite to the Infinite: Sharper Asymptotic Bounds on Norin's Conjecture via SAT}
\author{Markus Kirchweger, Tom\'a\v{s} Peitl, Bernardo Subercaseaux, Stefan Szeider}
\date{Source: arXiv:2511.08386v1 (CC BY 4.0)}
\begin{document}
\maketitle

\noindent\emph{Source and licence.} From the Finite to the Infinite: Sharper Asymptotic Bounds on Norin's Conjecture via SAT. Version arXiv:2511.08386v1 (CC BY 4.0), distributed under the Creative Commons Attribution 4.0 International licence, as recorded in the arXiv licence metadata for this exact version and shown on the abstract page https://arxiv.org/abs/2511.08386v1. Full rights notice: licenses/arxiv-2511.08386-CC-BY-4.0.txt.\par\smallskip

\noindent\textit{Editorial scope.} Counted unit: the Lemma whose source label is \texttt{None} in the article (source0.tex lines 9673--9679, source label \texttt{None}), delivered together with the article's own complete original proof of it (source0.tex lines 9680--9735, one \texttt{proof} environment of 56 lines). No mathematics is written, repaired or re-derived: statement and proof are the article's own text line for line. Required context retained uncounted: none. Every definition, display and label of those lines is retained unchanged. Omitted from this package and not redistributed: 1--9672, 9736--9821. Nothing inside a retained excerpt is omitted or altered: every retained body is delivered exactly as the article's lines are written, so no omission receipt of any kind accompanies this package. Citations: the retained excerpts carry no citation command at all, so no bibliography entry of the article is transcribed and no reference is redistributed. Reference adaptation: none was needed and none was made; every reference inside the delivered unit resolves inside it, so no unresolved reference or citation is produced. Printed slips kept literal: no printed word, symbol, hypothesis or number of the counted statement or of its proof was altered. Layout and class: the article's own class is \texttt{article} with packages including booktabs, algorithm, algorithmic, babel, xcolor, subcaption, nicefrac, url, fullpage, tikz, pgfplots, geometry, setspace, csquotes; the wrapper is the lane's pinned \texttt{amsart} editorial wrapper at 10pt on A4, which restates the theorem-like environments the retained text needs and adds the article's own macro definitions that the retained text needs (\texttt{\textbackslash leq}), with extra wrapper packages (bm, bbm, dsfont, xspace, cancel, mathtools). The delivered numbering is the unit's own. Assets: no retained span contains a figure environment, a graphics-inclusion command, a PDF-page inclusion, a table or a TikZ picture; no image file is copied into this package, the package declares no asset dependency and no image file is redistributed with it. Licence: version arXiv:2511.08386v1 of this article is distributed under the Creative Commons Attribution 4.0 International licence, as recorded in the arXiv licence metadata for this exact version and shown on the abstract page https://arxiv.org/abs/2511.08386v1. A full rights notice is delivered at licenses/arxiv-2511.08386-CC-BY-4.0.txt. Characters: every retained excerpt is pure ASCII, so the delivered engine, XeLaTeX with its default Latin Modern text font, reports no missing character.
\begin{lemma}
    Let $n \in \mathbb{N}$ be odd, let $c : E(Q_n) \to \{0,1\}$ be defined by:
    \[
        ((x_1, \dots, x_i, \dots, x_n),\, (x_1, \dots, 1 - x_i, \dots, x_n)) \mapsto \bigoplus_{j \neq i} x_j.
    \]
    The number of blocking antipodal pairs in $c$ is $2^{n-1}\left(1-O\left(\frac{1}{\sqrt{n}}\right)\right)$
\end{lemma}

\begin{proof}
Consider a path $u\rightarrow v \rightarrow w$, and suppose $u$ and $v$ differ at position $i$ and $v$ and $w$ at position $j$.
The color of the edge $uv$ is $c_{uv} = \bigoplus_{l\not\in \{i, j\}} u_l \; \oplus u_j$.
The color of the edge $vw$ is $c_{vw} = \bigoplus_{l\not\in \{i, j\}} v_l \; \oplus v_i$.
We have $c_{uv} \oplus c_{vw} = u_j \oplus v_i = u_j \oplus (1-u_i) = 1 \oplus u_j \oplus u_i$.
So, there is a color change at $v$ if and only if $u_i = u_j$, i.e., the two consecutive coordinates both flip $0$ to $1$ or both $1$ to $0$.

An antipodal geodesic from some vertex $v$ must flip all coordinates one by one.
If the numbers of $0$s and $1$s in $v$ differ by at most one, the flips can be interleaved in such a way that no color changes are needed.
If they differ by $2$, it can be done with one color change.
If they differ by at least $3$, two color changes are needed.

Let $\beta(v)$ be the difference between the number of $1$s and $0$s in $v$, namely $\beta(v) = \sum v_i - \sum (1-v_i) = 2\sum v_i - n = \sum v_i - \sum \bar{v}_i = -\beta(\bar v)$.
Represent an antipodal pair by the vertex which has more $0$s than $1$s, i.e., where $\beta(v) < 0$ (recall $n$ is odd).
We want to count the number of $v$ for which $\beta(v) \leq -3$.
We have $\beta(v) = 2\sum v_i - n \leq -3$ if and only if $\sum v_i \leq \frac{n-3}{2}$.
The number of such $v$ is
\[ g(n) = \sum_{i = 0}^{\frac{n-3}{2}} \binom{n}{i} .\]
%Some values of $g$ are $g(3) = 1, g(5) = 6, g(7) = 29, g(9) = 130$.

Write $n=2k+1$, let $h(k) = g(2k+1)$, and consider that

\begin{align*}
    2^{2k+1} = 2^n = \sum_{i=0}^{n} \binom{n}{i} = 
    2h(k) + \binom{2k+1}{k} + \binom{2k+1}{k+1} = 
    2\left(h(k) + \binom{2k+1}{k}\right),
\end{align*} 

and thus

\[ h(k) = 2^{2k} - \binom{2k+1}{k} .\]
%We aim to show for $k \geq 4$ that $h(k) > 2^{2k-1}$, i.e., equivalently that $\binom{2k+1}{k} < 2^{2k-1}$.
Plug in
\[ \binom{2k+1}{k} = \binom{2k}{k} + \binom{2k}{k-1} = \binom{2k}{k} \left(1 + \frac{k}{k+1}\right),\]
and apply the upper bound on the central binomial coefficient that arises from Stirling's approximation,
\( \binom{2k}{k} \leq \frac{2^{2k}}{\sqrt{\pi k}}\),
to get
\begin{align*}
\binom{2k+1}{k} \leq \left(1 + \frac{k}{k+1}\right)\frac{2^{2k}}{\sqrt{\pi k}} =
%O \left( 2^{2k} \frac{1}{\sqrt{k}} \right) =
O \left( 2^{n-1} \frac{1}{\sqrt{n}} \right).
%= 2^{2k}\left(\frac{2k+1}{(k+1)\sqrt{\pi k}}\right) < 2^{2k}\left(\frac{2k+1}{(k+1)\sqrt{3k}}\right).
\end{align*}
%We already know that $130 = h(4) > 2^{2\cdot 4 -1} = 128$, so let us assume $k \geq 5$.
%We need to show that for $k \geq 5$
%\[ \frac{2k+1}{(k+1)\sqrt{3k}} \leq \frac{1}{2},\]
%which is equivalent to
%[ 4(2k+1)^2 \leq (k+1)^2 3k,\]
%which is equivalent to
%\[ 3k^3 - 10k^2 - 13k - 4 \geq 0 .\]
%It can be shown that the LHS has exactly one positive root, which lies between $4$ and $5$, and is positive %thereafter, completing the proof.

%More generally, since
%\[ \lim_{n \to \infty} \frac{g(n)}{2^{n-1}} = \lim_{k \to \infty} \frac{h(k)}{2^{2k}} = \lim_{k \to \infty} 1 - %\frac{\binom{2k+1}{k}}{2^{2k}} \geq 1 - \lim_{k \to \infty} \frac{2k+1}{(k+1)\sqrt{3k}} = 1, \]
%it holds that for any $\varepsilon>0$ for sufficiently large $n$ at least $1-\varepsilon$ fraction of antipodal pairs are bad (require at least two switches). In fact, the proportion of good pairs is $O(\frac{1}{\sqrt{k}})$.
\end{proof}

\end{document}
